\chapter{Elliptic problems: elasticity and heat transfer}
\label{chap:elasticity}

\framepart{Outline}

This chapter sets up and solves the two model elliptic problems: linear
elastic equilibrium and thermal equilibrium. The two are treated together
because in \texttt{Cast3M} they are literally the same sequence of operators
with different keywords -- build a model, attach a material, assemble a
matrix, impose boundary conditions, solve, post-process. Once you have the
elastic chain in your fingers the thermal one costs five minutes.

We continue with the plate with a hole meshed in
chapter~\ref{chap:meshing}.

\section{Variational formulation of the elastic equilibrium}

\todoeq{variational formulation -- see the corresponding chapter of the
lecture notes}

\section{General settings}

The construction of a finite element solution starts with the definition of
the main modelling assumptions: dimension of the space, coordinate system,
plane strain or plane stress, Fourier representation, and so on.

\begin{castemcom}
The dimension of the working space and the fundamental assumptions of the
computation are chosen with the \op{opti} command. The two options shown here
correspond to plane strain, i.e. \op{defo[rmation] plan[e]}
\dindex{defo[rmation] plan[e]}{plane strain} in cartesian coordinates, and to
axial symmetry, i.e. \op{axis} \dindex{axis}{axial symmetry} in cylindrical
coordinates.
\end{castemcom}
\begin{castemlines}
\begin{verbatim}
opti mode defo plan;

opti mode axis;
\end{verbatim}
\end{castemlines}

\noindent The chosen spatial configuration automatically assigns the standard
names of the components of the different fields. \textbf{You do not get to
choose these names, and you will need them constantly} -- every \op{exco},
every \op{bloq}, every plot refers to a component by name.

\begin{itemize}
\item \textbf{plane strain} (\op{defo[rmation] plan[e]}). Displacement fields
  are defined in a cartesian coordinate system $(x,y,z)$ as
  \[ \vect{u}(x,y,z)=u_{x}(x,y)\,\vect{e}_{x}+u_{y}(x,y)\,\vect{e}_{y} \]
  and the component names are:

  \begin{center}
  \begin{tabular}{@{}l@{\hspace{10mm}}l@{}}
    \textbf{field} & \textbf{component name}\\[1mm]
    displacement & \texttt{UX, UY}\\
    forces       & \texttt{FX, FY}\\
    strain       & \texttt{EPXX, EPYY, EPXY, EPZZ}\\
    stress       & \texttt{SMXX, SMYY, SMXY, SMZZ}\\
  \end{tabular}
  \end{center}

  Note that \texttt{EPZZ} and \texttt{SMZZ} are present: in plane strain
  $\varepsilon_{zz}=0$ but $\sigma_{zz}=\nu(\sigma_{xx}+\sigma_{yy})\neq0$,
  and it is a standard mistake to forget it when computing an equivalent
  stress by hand.

\item \textbf{axial symmetry} (\op{axis}). Displacement fields are defined in
  a cylindrical coordinate system $(r,\theta,z)$ as
  \[ \vect{u}(r,\theta,z)=u_{r}(r,z)\,\vect{e}_{r}+u_{z}(r,z)\,\vect{e}_{z} \]
  and the component names become:

  \begin{center}
  \begin{tabular}{@{}l@{\hspace{10mm}}l@{}}
    \textbf{field} & \textbf{component name}\\[1mm]
    displacement & \texttt{UR, UZ}\\
    forces       & \texttt{FR, FZ}\\
    strain       & \texttt{EPRR, EPZZ, EPRZ, EPTT}\\
    stress       & \texttt{SMRR, SMZZ, SMRZ, SMTT}\\
  \end{tabular}
  \end{center}

  Here the hoop terms \texttt{EPTT}, \texttt{SMTT} carry the
  $\theta\theta$ component. In an axisymmetric computation the mesh must lie
  in the half-plane $r\geq0$, with $r$ as the first coordinate.
\end{itemize}

\noindent The complete list for every modelling option is in the notice of
\op{mode[le]} on the website~\cite{Cast3M:web}; the exact set also depends on
the element, so check it there rather than trusting this table for an element
type not used in this booklet.

\begin{optable}{General settings}
opti[on] dime n   & space dimension $n=1,2,3$ \\
opti[on] mode plan defo & cartesian coordinates, plane strain \\
opti[on] mode plan cont & cartesian coordinates, plane stress \\
opti[on] mode axis      & cylindrical coordinates, axisymmetric \\
opti[on] mode four[ier] n & cylindrical coordinates, Fourier expansion up to
                            $n$ harmonics \\
opti[on] mode trid[imensionnel] & cartesian coordinates, three-dimensional \\
\end{optable}

\section{Elasticity -- \texttt{Cast3M} problem setting}

\subsection{Boundary conditions and solution method}

The method of solution used by \texttt{Cast3M} is based on the Lagrangian
relaxation of the displacement boundary conditions, as explained for example
in~\cite{Bonnet:2007}. The variational principle
of the preceding section transforms into

\noindent which, after discretisation, leads to the linear
system~\cite{Bonnet:2007}

\begin{equation}
\begin{bmatrix} K & A^{T} \\ A & 0 \end{bmatrix}
\begin{bmatrix} U \\ \lambda \end{bmatrix}
=
\begin{bmatrix} F \\ U^{D} \end{bmatrix}
\label{eq:elas:lagrange}
\end{equation}

\noindent in which the stiffness matrix $[K]$ is defined with respect to
\emph{all} nodal displacements, with no distinction between values imposed by
the boundary conditions and the others. The imposed displacements enter
through the projection matrix $[A]$ and the imposed values $[U^{D}]$. The
traction boundary conditions and the body forces define the vector $[F]$.

The practical consequence of \eqref{eq:elas:lagrange} is the one that
surprises newcomers: \textbf{imposing a boundary condition adds unknowns to
the system} -- the Lagrange multipliers $\lambda$ -- rather than removing
them. Those multipliers are the reaction tractions, which is why they can be
recovered afterwards with \op{reac}.

\begin{castemcom}
The \op{mode[le]} \dindex{mode[le]}{model} operator associates to a given mesh
a finite element, i.e. a specific shape function in the variational principle
(here linear functions by default, the \texttt{QUA4} element), a mechanical
problem setting \op{mecanique}, and a linear isotropic elastic constitutive
law. Other problem settings are \op{thermique}, \op{liquide}, \op{darcy},
\ldots

The choice of the material behaviour determines the number and the names of
the parameters that must be supplied next through the \op{mate[riaux]}
\dindex{mate[riaux]}{material} operator, as well as the structure of the
internal variables of the model.

The stiffness matrix $[K]$ is then computed with \op{rigi[dite]}
\dindex{rigi[dite]}{stiffness} from the elastic information provided before.
\end{castemcom}
\begin{castemlines}
\begin{verbatim}
mod = dom mode mecanique
          elastique isotrope;

mat = mod mate 'YOUN' 2.e11 'NU' 0.3;

KK = rigi mod mat;
\end{verbatim}
\end{castemlines}

\begin{castemcom}
To construct the projection matrices $[A]$, which reduce the displacement
vector to the part of the boundary where $[U^{D}]$ is applied, we use
\op{bloq[uer]} \dindex{bloq[uer]}{block a degree of freedom}. The
corresponding right-hand side is then built from $[A]$ and $[U^{D}]$ with
\op{depi} \dindex{depi}{imposed displacement} (fr. \emph{déplacement imposé}).

Note that \op{bloq[uer]} and \op{depi} do not merely build $[A]$ and
$[U^{D}]$: they also create the new degrees of freedom corresponding to the
Lagrange multipliers. Mechanically these multipliers are the tractions $[T]$
on that part of the boundary, unknown before the system is solved. You can
see them by listing the object: \op{list Ad34;}.

Here \op{d12} and \op{d41} carry the symmetry conditions of the quarter
plate ($u_y=0$ on the lower edge, $u_x=0$ on the left edge), while the block
on \op{d34} is the loading: the upper edge is given the imposed displacement
$u_y=0.1$.
\end{castemcom}
\begin{castemlines}
\begin{verbatim}
Ad12 = bloq uy d12;
Ad41 = bloq ux d41;
Ad34 = bloq uy d34;

ud34 = depi Ad34 0.1;

AA = Ad12 et Ad41 et Ad34;
\end{verbatim}
\end{castemlines}

\begin{castemcom}
The system is finally solved by calling \op{reso[udre]}
\dindex{reso[udre]}{solve}. Only the non-zero components of the right-hand
side have to be specified -- here only \op{ud34}; the rest is completed with
zeros automatically.
\end{castemcom}
\begin{castemlines}
\begin{verbatim}
u = reso (KK et AA) ud34;
\end{verbatim}
\end{castemlines}

\begin{castemcom}
Listing the components of the displacement field shows the expected
\texttt{UX} and \texttt{UY}, and also a new component \texttt{LX}: the
Lagrange multiplier, which carries $[T]$.

To obtain the underlying traction vector we use \op{reac[tion]}
\dindex{reac[tion]}{reaction forces}, which turns the multipliers back into
nodal forces on the blocked boundary. Summing them is the standard way to get
the resultant force on a boundary -- for instance to plot a force/displacement
curve.
\end{castemcom}
\begin{castemlines}
\begin{verbatim}
list (extr u 'COMP');

td34 = reac u Ad34;

* resultant on d34
ftot = resu (exco td34 'FY');
\end{verbatim}
\end{castemlines}

\subsection{Have the boundary conditions removed every rigid body motion?}
\label{sec:elas:ense}

\cindex{rigid body motion}\cindex{under-restrained structure}%
\cindex{singular stiffness matrix}%
Before solving anything, it is worth asking whether the blocked degrees of
freedom actually hold the structure. A body that can still translate or
rotate as a whole has a singular stiffness matrix, and the symptoms are
unmistakable once you know them: the solver fails, or it returns
displacements of absurd magnitude with a sensible-looking stress field
superposed on them.

In two dimensions a body has three rigid body motions -- two translations and
one rotation -- and in three dimensions six. Every one of them must be
removed by the boundary conditions. Counting them by hand on the quarter
plate is easy; on a real assembly of several parts, with symmetry conditions
and contact, it is not, and this is where most "the solver does not converge"
problems actually live.

\texttt{Cast3M} will answer the question directly.

\begin{castemcom}
\op{ense[mble]} \dindex{ense[mble]}{rigid body modes of a stiffness matrix}
takes an assembled stiffness -- \emph{including its blocking matrices} -- and
returns a \op{solution} object containing the \emph{mouvements d'ensemble},
the rigid body modes that the matrix still permits.

\medskip
\textbf{The test is the number of modes it finds.} None means the structure
is properly restrained and you can solve. One or more means it is not, and
each mode it returns \emph{is} a motion your \op{bloq} conditions failed to
suppress -- so you do not have to guess which one is missing: plot it.

\op{tire} \dindex{tire[r]}{extract a mode from a solution object} pulls out
mode number \texttt{n} as a displacement field, which is then drawn with
\op{defo} like any mode shape.
\end{castemcom}
\begin{castemlines}
\begin{verbatim}
* KK is the stiffness, AA the
* assembled blocking matrices
mmm = ense (KK et AA);

* how many rigid modes remain?
nrig = dime mmm;
mess 'rigid body modes left: ' nrig;

* ... and what are they?
si (nrig > 0);
  repeter brig nrig;
    i  = &brig;
    di = tire mmm 'DEPL' 'RANG' i;
    trac (defo dom di 1. rouge);
  fin brig;
finsi;
\end{verbatim}
\end{castemlines}

\noindent The complement of \op{bloq} is therefore \op{ense}: one imposes the
conditions, the other checks that they were enough. Running it once on every
new model costs nothing and saves a great deal of time.

\medskip
Two situations where you will want it even though nothing has gone wrong:

\begin{itemize}
\item \textbf{a free-free structure, deliberately.} A modal analysis of an
  unrestrained body -- an aircraft, a satellite panel, a specimen hanging on
  soft springs -- is supposed to have rigid body modes, and
  \op{vibr[ation]} will return them as modes at (numerically) zero frequency.
  Knowing from \op{ense} how many to expect is exactly how you tell a genuine
  zero-frequency mode from a mesh that has fallen apart;
\item \textbf{before a buckling computation.} \op{flambage} works on a
  stiffness that must be non-singular, so an under-restrained model gives
  meaningless critical loads rather than an error. Chapter~\ref{chap:eigen}
  returns to both of these.
\end{itemize}

\noindent A related use of the word: \op{rela 'ENSE'}, in
section~\ref{sec:elas:rela}, imposes that a component be \emph{equal} over a
whole boundary -- a rigid platen. The two are unconnected beyond the shared
French word \emph{ensemble}, "as a whole".

\medskip
\noindent\textbf{A note on the name.} The operator is \op{ense[mble]}, not
\op{enve}. \op{enve[loppe]} is the meshing operator of
section~\ref{sec:mesh:3d} that extracts the boundary surface of a volume
(and, in a second form, the envelope of a set of oscillator spectra). The two
are easy to confuse from memory and do entirely different things.

\begin{optable}{Plots}
trac[er] & plots different objects \\
trac[er] model field & plots a field defined on elements \\
trac[er] geometry field & plots a field defined at nodes \\
trac[er] geometry & plots a mesh \\
dess[iner] evolution & plots the graph of a function \\
evol[ution] & graph of a function defined by lists of numbers \\
evol[ution] 'CHPO' \ldots & extracts the evolution of a field along a path \\
\end{optable}

\noindent The results can be post-processed in a number of ways, as shown in
the next examples.

\begin{castemcom}
To plot the norm of the displacement vector
\[ \norm{u}=\sqrt{u_{x}^{2}+u_{y}^{2}} \]
we extract the components with \op{exco} \dindex{exco}{extract a component}
(fr. \emph{extraire composante}), combine them algebraically, and plot with
\op{trac[er]}.

The second argument of \op{exco} is the component to take, the third the name
to give it in the result: renaming to \texttt{'SCAL'} is what lets the two
components be multiplied together afterwards -- fields are only combined
component by component, by name.

As displacement fields are defined by their nodal values, the underlying mesh
must be given for the plot.
\end{castemcom}
\begin{castemlines}
\begin{verbatim}
uux = exco u 'UX' 'SCAL';
uuy = exco u 'UY' 'SCAL';

trace uux dom;

unorm = (((uux*uux) + (uuy*uuy))**0.5);
trace unorm dom;
\end{verbatim}
\end{castemlines}

\begin{castemcom}
Strains and stresses are computed with \op{epsi[lon]}
\dindex{epsi[lon]}{strain} and \op{sigm[a]} \dindex{sigm[a]}{stress}; the
operations correspond to the matrix computations
\begin{align}
\varepsilon &= [B][U] \\
\sigma      &= [C][B][U]
\end{align}
with $[C]$ the elasticity matrix. Note that \op{sigm} needs the material in
addition to the model, whereas \op{epsi} does not -- strain is a purely
kinematic quantity.

Both results are \op{mchaml} objects, defined at the Gauss points. The names
of their components can be listed with \op{extr[aire]}.
\end{castemcom}
\begin{castemlines}
\begin{verbatim}
eps = epsi mod u;
sig = sigm mod mat u;

list (extr eps 'COMP');
list (extr sig 'COMP');
\end{verbatim}
\end{castemlines}

\begin{castemcom}
Another operator corresponds to $[B^{T}]$ and computes the nodal forces
associated with a given stress field:
\begin{equation}
[F]=[B^{T}]\sigma
\end{equation}
which is the discrete counterpart of $\ddiv\sigma$ in a continuum
formulation. \op{bsig[ma]} \dindex{bsig[ma]}{nodal forces from a stress field}

A good sanity check on any converged solution: these nodal forces must vanish
everywhere except at the boundary, where they equal the reactions. If they do
not, the solution is not in equilibrium.
\end{castemcom}
\begin{castemlines}
\begin{verbatim}
f = bsigm mod sig;

trace dom (exco f 'FX');
trace dom (exco f 'FY');
\end{verbatim}
\end{castemlines}

\begin{optable}{Strains and stresses}
epsi[lon] & computes the strain tensor \\
sigm[a]   & computes the stress tensor \\
bsig[ma]  & computes the nodal forces corresponding to a stress tensor \\
tres[ca]  & computes the Tresca norm of the stress tensor \\
vmis[es]  & computes the von Mises norm,
            $\sqrt{\tfrac{3}{2}\,\dev\sigma\bcddot\dev\sigma}$ \\
prin[cipales] & computes the principal values of the stress tensor \\
@plotpri  & plots the principal values (user-defined procedure) \\
\end{optable}

\begin{castemcom}
To plot the components of a strain or stress tensor we apply \op{trac[er]}
directly and adjoin the \emph{model}, not the mesh, so that the position of
the Gauss points where the values are defined is known.

This is the practical difference between a \op{chpoint} and a \op{mchaml}:
\op{trace \textit{mesh} \textit{field}} for the first,
\op{trace \textit{model} \textit{field}} for the second.
\end{castemcom}
\begin{castemlines}
\begin{verbatim}
trace mod sig;
\end{verbatim}
\end{castemlines}

\begin{castemcom}
Different scalar measures derived from the stress tensor are computed
directly: the Tresca and von Mises norms with \op{tres[ca]} and
\op{vmis[es]} \dindex{vmis[es]}{Mises stress norm}, the principal values and
their directions with \op{prin[cipales]}.

\op{@plotpri} is a user-defined procedure which plots the principal stresses
interactively. Procedures whose names start with \op{@} are user
contributions rather than standard operators -- see
chapter~\ref{chap:importexport} on how to load them.
\end{castemcom}
\begin{castemlines}
\begin{verbatim}
sigvm = vmis mod sig;
sigtr = tres mod sig;
sigpr = prin mod sig;

@plotpri sig mod;
\end{verbatim}
\end{castemlines}

\begin{castemcom}
It is often more informative to display the evolution of a stress component
along a path than over the whole domain. Here we plot $\sigma_{yy}$ along the
border of the circular hole, where the stress concentration is.

We first list the components of the stress field. We then interpolate the
$\sigma_{yy}$ component -- \texttt{'SMYY'} -- from the Gauss points to the
nodes using \op{chan[ger]} with the \texttt{'CHPO'} option
\dindex{chan[ger] 'CHPO'}{change a field into nodal values}: a path extraction
needs nodal values. Finally \op{evol[ution]} with the \texttt{'CHPO'} option
extracts the values along the path \op{(arc23 et arc34)} and \op{dess[iner]}
draws the graph.
\sindex{evol[ution], 'CHPO', along a path}
\end{castemcom}
\begin{castemlines}
\begin{verbatim}
list (extr sig 'COMP');

sigyy = chan 'CHPO' mod
             (exco sig 'SMYY');

ev1 = evol 'CHPO' sigyy
           (arc23 et arc34);
dess ev1;
\end{verbatim}
\end{castemlines}

\noindent For the plate with a hole in tension this is the computation that
should return the classical stress concentration factor $K_t=3$ for an
infinite plate; comparing your value with 3 is the natural first validation of
the whole chain, and the discrepancy tells you how much the finite width of
the plate matters.

\subsection{Other applied forces: pressure, weight, \ldots}

The load in the previous example was an imposed displacement only, so the
exterior load vector was $[F]=[0]$ and the program completed it with zeros.
The next examples show how surface tractions and body forces are imposed.

\begin{optable}{Applied forces}
forc[e]    & distributes a given force over a set of nodes \\
pres[sion] & computes the nodal force distribution for a given pressure \\
coor[donnees] & computes the coordinate fields of a given mesh \\
mass[e]    & computes the mass matrix \\
manu 'CHPO' & constructs a nodal field by hand \\
\end{optable}

\begin{optable}{Creation of stiffness-type matrices}
rigi[dite]     & stiffness matrix \\
mass[e]        & mass matrix \\
cond[uctivite] & conductivity matrix \\
capa[cite]     & capacity matrix \\
conv[ection]   & stiffness corresponding to a convection boundary condition \\
\end{optable}

\subsection{Surface traction and applied pressure}

The main operators for constructing applied forces are \op{forc[e]}
\sindex{forc[e]} and \op{pres[sion]} \sindex{pres[sion]}.

\begin{castemcom}
The \op{coor} \dindex{coor[donnees]}{coordinate fields} operator builds the
coordinate fields of a mesh. Here we recover $x$ and $y$ at each point of the
line \op{d23}.

\op{pres} computes the nodal force field corresponding to the pressure
distribution $p(x,y)=3x^{2}$. A model is needed because the computation
integrates the pressure against the shape functions of the underlying
elements (see~\cite{Bonnet:2007}) -- a pressure is not simply divided equally
between nodes. Note
that this model must be built \emph{on the loaded boundary} \op{d23}, not on
the surface \op{dom}: it is the line elements that are integrated over. This
is a frequent source of errors.

Finally the normal pressure distribution \op{fnor} is turned into a tangential
traction \op{ftan} by exchanging the components of the vector field and
changing one sign, i.e. by rotating it by $90^{\circ}$.
\end{castemcom}
\begin{castemlines}
\begin{verbatim}
* the model must live on the loaded
* line, not on the surface
mod23 = d23 mode mecanique
            elastique isotrope;

xx yy = coor d23;

fnor = pres mass mod23 (3.* (xx * xx));

ftan = (exco fnor 'FY' 'FX') et
       (-1. * (exco fnor 'FX' 'FY'));
\end{verbatim}
\end{castemlines}

\noindent To impose a Hertzian pressure distribution, for example a parabola
centred on the left end of the loaded edge,
\[
p(x)=\begin{cases}
p_{max}\,(a_{p}^{2}-x^{2}) & \text{if } \abs{x}<a_{p}\\
0 & \text{if } \abs{x}\ge a_{p}
\end{cases}
\]
we proceed in the same way, with one extra ingredient.

\begin{castemcom}
\op{manu[el]} \dindex{manu[el]}{create a field by hand} constructs a constant
nodal field, which is what lets a constant appear in an otherwise
field-valued formula: $a_p$ has to become a field before $a_p^2-x^2$ can be
formed.

\op{masq[uer]} \dindex{masq[uer]}{mask, characteristic function} builds the
characteristic function of the condition
$x<a_{p}$ -- a field equal to 1 where the condition holds and 0 elsewhere.
Multiplying by it is how one writes a piecewise definition in
\textsc{gibiane}; there is no \op{si} test inside a field expression. Here
the edge runs from $x=0$, so $x<a_p$ and $\abs{x}<a_p$ select the same
points; on an edge straddling the origin you would need the product of two
masks, as in the inclusion example below.

The final nodal force distribution is assembled by multiplying the two.
\end{castemcom}
\begin{castemlines}
\begin{verbatim}
xx = coor 1 d34;

ap   = ldom / 3.;
pmax = 1.e8;

one = manu 'CHPO' d34 1 'SCAL' 1.;

px  = pmax * (((ap * one) ** 2)
              - (xx ** 2));

mask = (xx masq 'INFERIEUR' ap);
px   = px * mask;

mod34 = d34 mode mecanique
            elastique isotrope;
fp = pres mass mod34 px;
\end{verbatim}
\end{castemlines}

\subsection{Body forces}

\begin{castemcom}
To build a body force taking the density of the body into account, the
density \texttt{RHO} is added to the material and the mass matrix is computed
with \op{mass[e]}.

If the gravitational acceleration is defined as the nodal field \op{g}, the
body force is simply the product of the mass matrix with it: $[F]=[M]\{g\}$.
That is the discrete form of $\int_\Omega \rho\, \vect{g}\cdot\vect{v}$, and
it is the correct way to apply gravity -- distributing the total weight
equally over the nodes is not.
\end{castemcom}
\begin{castemlines}
\begin{verbatim}
mat = mod mate 'YOUN' 2.e11 'NU' 0.3
               'RHO' 2000.;

MM = mass mod mat;

g = manu 'CHPO' dom 1 'UY' -9.81;

fvol = MM * g;
\end{verbatim}
\end{castemlines}

\subsection{Thermal expansion}

\cindex{thermal expansion}\cindex{thermoelasticity}%
\cindex{dilatation coefficient}%
The thermal expansion of a heated body is a loading like any other, and it is
worth treating here because it is the commonest way a thermal computation
feeds back into a mechanical one. It follows directly from the thermoelastic
constitutive law
\begin{equation}
  \tens{\sigma} = \tens{C} : \left( \tens{\varepsilon} - \tens{a}\,\theta \right)
\label{eq:elas:thermoelastic}
\end{equation}
where $\tens{C}$ and $\tens{a}$ are the fourth-order tensor of elastic moduli
and the second-order tensor of thermal dilatation, and $\theta$ is the field of
temperature difference between the initial and the actual configuration. In the
elastic problem setting $\theta$ is known: it plays the role of a loading.

For an isotropic material this simplifies to
\begin{equation}
  \tens{\sigma} = \lambda \tr (\tens{\varepsilon} - \alpha \theta \tens{I}) \tens{I}
   + 2 \mu (\tens{\varepsilon} - \alpha \theta \tens{I})
\end{equation}
and, in the absence of body forces, the balance equation
$\ddiv \tens{\sigma} = 0$ becomes
\begin{equation}
   \ddiv \left[ \tens{C} : \tens{\varepsilon} \right]
   = \ddiv \left[ \tens{C} : \left( \tens{a}\,\theta \right) \right]
\end{equation}

\noindent The right-hand side is the point: a temperature field acts exactly
like a distributed body force, generated by the incompatibility of the thermal
expansion. That is why a uniform temperature rise on an unconstrained body
produces no stress at all, while the same rise on a constrained or
inhomogeneous body does.

These steps map one to one onto the \textsc{gibiane} program.

\begin{castemcom}
The thermal expansion coefficient must be declared in the material data. For
an isotropic elastic material the component name is \texttt{'ALPH'}.

\medskip
In this example the temperature field $\theta$ is \op{delta\_th}. The operator
\op{thet[a]} \dindex{thet[a]}{thermal stress field} computes, from the
material data and the temperature field, the thermal stress field
\[ \op{sig\_th} = \tens{C} : \left( \alpha\,\theta\,\tens{I} \right) \]
The equivalent body forces follow from \op{bsig[ma]}, whose role was described
earlier in this chapter.

The solution is then obtained with the standard \op{reso[udre]}, from the
stiffness matrix \op{KK}, the matrices of the imposed displacement conditions
\op{AA}, and the body force field \op{f\_th}.

\medskip
Note that \op{thet} measures the temperature from a reference value
\texttt{'TALP'}, which defaults to zero: \op{delta\_th} is a temperature
\emph{difference}, not an absolute temperature.
\end{castemcom}
\begin{castemlines}
\begin{verbatim}
mod = dom mode mecanique
          elastique isotrope;
mat = mod mate 'YOUN' 2.e11 'NU' 0.3
               'ALPH' 1.e-5;

* ... boundary conditions as before,
*     giving KK and AA ...

delta_th = manu 'CHPO' dom 1 'T' 1000.;

sig_th = thet mod mat delta_th;
f_th   = bsigm mod sig_th;

u = reso (KK et AA) f_th;
\end{verbatim}
\end{castemlines}

\noindent The stress actually carried by the body is \emph{not} \op{sig\_th}:
it is the stress computed from the resulting displacement minus the thermal
part, i.e. \op{(sigm mod mat u) - sig\_th}. Forgetting that second term is the
classic error in a thermoelastic computation, and it shows up as a non-zero
stress in a body that is free to expand.

\subsection{Other constraints on displacements}
\label{sec:elas:rela}

Blocking a degree of freedom to a given value is not the only kind of boundary
condition. \cindex{rigid platen}\cindex{linear constraint between unknowns}%
A frequent requirement is that a part of the boundary stay flat or
move as a rigid whole -- a loading platen, a symmetry condition on a surface
that is free to translate -- i.e. that all the displacement components of the
points of a line $\ell$ be equal to each other:
\begin{equation}
   \vect{u}(\vect{x}) = \text{constant}
   \qquad \forall \vect{x} \in \ell \subset \partial \Omega
\end{equation}
The constant itself is left free, and is fixed afterwards -- for instance by
imposing the displacement at one particular point of the line, as in the
example below.

\dindex{rela[tion]}{constraint on the displacement field}

\begin{castemcom}
The operator \op{rela[tion]} with the option \texttt{'ENSE'}
(fr. \emph{ensemble}) builds a stiffness matrix -- again a set of Lagrange
multipliers -- imposing that the \texttt{UY} component of the displacement
field be the same at every node of the line \op{d34}.

To fix the value of that unknown constant we impose the displacement at the
single point \op{p3}. The rest is as before: assemble all the matrices and
solve with \op{reso[udre]}.

\medskip
This is what you want for a rigid platen pressing on the plate: the surface
stays flat, but the traction distribution under it is computed rather than
assumed -- which is exactly the difference between this and \op{bloq}-ing the
whole line to a fixed value.
\end{castemcom}
\begin{castemlines}
\begin{verbatim}
Ad34 = rela 'ENSE' 'UY' d34;

Ap3 = bloq 'UY' p3;
up3 = depi Ap3 0.3;

AA = Ad12 et Ad41 et Ad34 et Ap3;

u = reso (KK et AA) up3;
\end{verbatim}
\end{castemlines}

\noindent \op{rela[tion]} imposes many other kinds of constraint, equalities
and inequalities alike: rigid body motion (\texttt{'CORI'}), sliding relations
between two meshes (\texttt{'GLIS'}), tying mid-side nodes to corner nodes
(\texttt{'MILI'}), coupling a pipe surface to its centre line
(\texttt{'TUYA'}), and so on. See its notice for the full list.

\begin{optable}{Conditions on displacements}
bloq[uer] & blocks a degree of freedom to an imposed value \\
rela[tion] & stiffness matrix for a linear relation between degrees of
             freedom \\
rela[tion] 'ENSE' & imposes one component to be equal over a whole mesh \\
rela[tion] 'CORI' & imposes a rigid body motion \\
depi      & right-hand side associated with an imposed value \\
reac[tion] & recovers the reaction associated with a constraint \\
\end{optable}

\subsection{Spatially varying material coefficients}

We have seen that spatially varying boundary tractions are easily defined with
\op{coor[donnees]} and \op{masq[uer]}. The same construction gives varying
material parameters. The example below places a square inclusion in a plate --
a two-material composite.

\begin{castemcom}
The example runs on its own square mesh \op{plaq} rather than on the plate
with the hole, since an inclusion spanning $0.2\le x,y\le 0.8$ would otherwise
overlap the hole.

The characteristic field of the inclusion, \op{D\_CAR}, is built as the
product of four masks and takes the value 1 inside the square and 0 outside;
\op{M\_CAR} is its complement.

The distributions of Young's modulus and Poisson's ratio are then
interpolations between the two materials weighted by these two fields.

The next step is the one to pay attention to. Material parameters live at the
Gauss points of the elements, not at the nodes, so the nodal fields must be
converted with \op{chan[ger]}
\dindex{chan[ger] 'CHAM'}{change from nodal to element field} using the
keywords \texttt{'CHAM'} (fr. \emph{champ}) and \texttt{'RIGIDITE'}, which
define the intended use and the header of the new field.
\dindex{chpo}{change from nodal to element field}%
\dindex{mchml}{change from nodal to element field} The component name
also has to be changed from \texttt{'SCAL'} to \texttt{'YOUN'} and
\texttt{'NU'} with \op{nomc} \dindex{nomc}{rename a component}, because
\op{mate} identifies its arguments by component name.

Material and stiffness matrix are then created as usual, the only difference
being that the arguments are fields rather than numbers.
\end{castemcom}
\begin{castemlines}
\begin{verbatim}
* a plain square of side 1, not the
* plate with the hole: the inclusion
* below would overlap the hole
opti dime 2 elem qua4;
c1 = 0. 0.;  c2 = 1. 0.;
c3 = 1. 1.;  c4 = 0. 1.;
plaq = dall (droi 20 c1 c2) (droi 20 c2 c3)
            (droi 20 c3 c4) (droi 20 c4 c1);

mod = plaq mode mecanique
           elastique isotrope;

X_min = 0.2; X_max = 0.8;
Y_min = 0.2; Y_max = 0.8;

xx yy = coor plaq;

D_CAR = xx masq 'SUPERIEUR' X_min;
D_CAR = D_CAR *
        (xx masq 'INFERIEUR' X_max);
D_CAR = D_CAR *
        (yy masq 'SUPERIEUR' Y_min);
D_CAR = D_CAR *
        (yy masq 'INFERIEUR' Y_max);

one   = manu 'CHPO' plaq 1 'SCAL' 1.;
M_CAR = one - D_CAR;

* aluminium and copper
YG_al = 69.e9;  NU_al = 0.33;
YG_co = 117.e9; NU_co = 0.35;

YG_r = (YG_co * M_CAR)
       + (YG_al * D_CAR);
NU_r = (NU_co * M_CAR)
       + (NU_al * D_CAR);

yg_c = chan 'CHAM'
       (nomc 'YOUN' YG_r) mod 'RIGIDITE';
nu_c = chan 'CHAM'
       (nomc 'NU'  NU_r) mod 'RIGIDITE';

trace mod yg_c;

mat = mate mod 'YOUN' yg_c 'NU' nu_c;
rig = rigi mod mat;
\end{verbatim}
\end{castemlines}

\noindent Two remarks on this example. First, the moduli in the original
version of these notes were identical for the two materials
($33\,$GPa and $\nu=0.32$ for both aluminium and copper), so the computation
returned a homogeneous plate; the values above are the usual ones and actually
produce an inclusion. Second, a continuously varying distribution is obtained
in exactly the same way, replacing the masks by any expression in \op{xx} and
\op{yy}.

\begin{optable}{Manipulation of fields}
chan[ger] & changes the field type by interpolation \\
manu[el]  & creates new objects \\
nomc      & changes the name of a component \\
exco      & extracts a component, optionally renaming it \\
masq[uer] & characteristic function of a condition \\
\end{optable}

\section{Thermal equilibrium -- \texttt{Cast3M} problem setting}
\label{sec:elas:thermal}

The steady thermal equilibrium problem is
\begin{equation}
  \ddiv \left( k \, \grad \theta \right) + r = 0
  \qquad \text{in } \Omega
\label{eq:thermal:equilibrium}
\end{equation}
with $k$ the conductivity and $r$ a volume heat source, completed by imposed
temperature, imposed flux or convection conditions on the boundary. It has
exactly the structure of the elastic problem, and is solved by exactly the
same sequence of operators.

\begin{castemcom}
The finite element model and the material properties are declared exactly as
in the elastic case, with \op{mode[le]} and \op{mate[riaux]}; only the
keywords change.

The resolution of~\eqref{eq:thermal:equilibrium} depends only on the
conductivity \texttt{K}. If the
density \texttt{RHO} and the specific heat \texttt{C} are also given, the
capacity matrix can be computed as well, which is what the transient
computation of chapter~\ref{chap:heat} will need.

\op{cond[uctivite]} and \op{capa[cite]} play exactly the roles that
\op{rigi[dite]} and \op{mass[e]} play in elasticity -- same objects, same
algebra, same solver.
\sindex{thermique} \sindex{conductivite} \sindex{capacite}
\end{castemcom}
\begin{castemlines}
\begin{verbatim}
mod = dom mode thermique isotrope;

mat = mod mate 'K' 1. 'RHO' 1. 'C' 1.;

KK = cond mod mat;
CC = capa mod mat;
\end{verbatim}
\end{castemlines}

\begin{castemcom}
An imposed temperature on the boundary uses the same Lagrange multiplier
technique as in elasticity: \op{bloq[uer]} builds the projection matrix and
\op{depi} the associated right-hand side.

The name of the degree of freedom is \texttt{'T'}, which is also the name of
the component of the temperature field.
\end{castemcom}
\begin{castemlines}
\begin{verbatim}
Ad23 = bloq 'T' d23;
AA   = Ad23;

Fd23 = depi Ad23 4.;
\end{verbatim}
\end{castemlines}

\begin{castemcom}
The nodal ``forces'' equivalent to a boundary heat flux are obtained with
\op{flux} \dindex{flux}{boundary heat flux}. As with \op{pres}, the model is
needed because the flux is integrated against the shape functions.
\end{castemcom}
\begin{castemlines}
\begin{verbatim}
Fd34 = flux mod 2. d34;
\end{verbatim}
\end{castemlines}

\begin{castemcom}
The convection boundary condition
\[ -k\derp{\theta}{n}=h\,(\theta-\theta_{\infty}) \]
is mathematically of a different type: it involves the unknown $\theta$ on the
boundary itself. It therefore contributes \emph{both} an extra conductivity
matrix \op{KK\_hol}, defined by the exchange coefficient $h$, \emph{and} a
nodal flux \op{Fhole} defined by $h$ and the environment temperature
$\theta_{\infty}$. Both must be added to the system.
\sindex{convection, boundary condition}

Radiation conditions of the type
$-k\derp{\theta}{n}=\varepsilon_{r}\sigma_{\!SB}(\theta^{4}-\theta_{\infty}^{4})$,
with $\sigma_{\!SB}$ the Stefan--Boltzmann constant and $\varepsilon_{r}$ the
emissivity, are
nonlinear and cannot be handled in this linear setting; a first approximation
is a temperature-dependent exchange coefficient $h=h(\theta)$, and the proper
treatment goes through \op{pasapas}.
\end{castemcom}
\begin{castemlines}
\begin{verbatim}
mod_hol = (arc23 et arc34)
          mode convection;
mat_hol = mod_hol mate 'H' 1.;

KK_hol = cond mod_hol mat_hol;
Fhole  = conv mod_hol mat_hol
              'T' 200.;
\end{verbatim}
\end{castemlines}

\begin{castemcom}
A distributed heat source inside the body is created with \op{sour[ce]}
\dindex{sour[ce]}{volume heat source}, given its value, the domain it applies
to and the model.
\end{castemcom}
\begin{castemlines}
\begin{verbatim}
Fsource = sour mod 100. dom;
\end{verbatim}
\end{castemlines}

\begin{castemcom}
The temperature field is finally obtained by solving the assembled system with
\op{reso[udre]}. Note that \emph{both} the conductivity matrix \op{KK} of
the domain and the convection matrix \op{KK\_hol} of the boundary condition
appear in the stiffness term -- that is what distinguishes a convection
condition from an imposed flux, which contributes only to the right-hand side.

All the right-hand side contributions are concatenated with \op{et}.
\end{castemcom}
\begin{castemlines}
\begin{verbatim}
KTOT = KK et KK_hol et AA;
FTOT = Fd23 et Fd34 et Fhole
            et Fsource;

TT = reso KTOT FTOT;

trace TT dom;
\end{verbatim}
\end{castemlines}

\noindent Varying boundary temperatures, fluxes or convection coefficients,
and volume heat sources, are imposed with the same \op{coor}/\op{masq}
constructions as the varying pressure above.

\begin{optable}{Thermal problem setting}
mode[le] \ldots thermique & thermal finite element model \\
mode[le] \ldots convection & model carrying a convection boundary condition \\
cond[uctivite] & conductivity matrix \\
capa[cite] & capacity matrix \\
flux & nodal flux equivalent to a boundary heat flux \\
conv[ection] & nodal flux of a convection condition \\
sour[ce] & nodal flux of a volume heat source \\
\end{optable}

\section{An unexpected use of the thermal solver}
\label{sec:elas:torsion}

\cindex{Laplace equation}\cindex{thermal analogy}%
\cindex{Saint-Venant torsion}\cindex{warping function}%
\cindex{torsional rigidity}\cindex{torsion centre}%
\cindex{Timoshenko beam}\cindex{shear correction coefficient}%
\cindex{beam cross-section properties}%
Before leaving the thermal operators, one idiom worth knowing, because it is
the kind of thing that distinguishes someone who knows a finite element code
from someone who only knows one physics in it: \textbf{the thermal solver can
be used to solve any Laplace problem}, whatever it means physically.

The classical example is the \emph{Saint-Venant torsion} of a prismatic beam.
When a bar of arbitrary cross-section is twisted, plane sections do not stay
plane -- they warp out of plane by a function $\varphi(x,y)$ which satisfies,
over the cross-section $S$,
\begin{equation}
  \Delta \varphi = 0 \quad \text{in } S,
  \qquad
  \derp{\varphi}{n} = y\,n_{x} - x\,n_{y} \quad \text{on } \partial S
\label{eq:elas:warping}
\end{equation}
and from which the torsional rigidity $J$ follows by integration over the
section. That is a Laplace equation with a prescribed normal derivative --
mathematically identical to steady heat conduction with an imposed boundary
flux, which is exactly what the previous section solved.

So the cross-section is meshed as a 2D domain, given a \emph{thermal} model
with unit conductivity, and the warping function is computed as if it were a
temperature.

\begin{castemcom}
The model is \op{thermique} even though the problem is mechanical. The
conductivity is set to 1, so that \op{cond} assembles exactly the discrete
Laplacian.

\op{flux} imposes the boundary term of~\eqref{eq:elas:warping}. A pure
Neumann problem is defined only up to a constant, so one point must be fixed
-- any point will do, and the choice shifts $\varphi$ without changing $J$.

The torsional rigidity is then an integral over the section, obtained with
\op{intg} \dindex{intg[rer]}{integrate a field over a domain}.
\end{castemcom}
\begin{castemlines}
\begin{verbatim}
opti dime 2 elem tri3;

* sect = mesh of the cross-section
modt = sect mode thermique isotrope;
matt = modt mate 'K' 1.;

KK = cond modt matt;

* fix the constant at one point
p0  = sect poin proc (0. 0.);
bl0 = bloq 'T' p0;
f0  = depi bl0 0.;

* ... boundary flux, built from the
* normal to the contour (see text)

phi = reso (KK et bl0) (fbord et f0);

trace phi sect;
\end{verbatim}
\end{castemlines}

\sindex{"@Torsion (procedure)}\cindex{Maxwell--Betti theorem}%
\noindent This is the basis of the \texttt{@Torsion} procedure presented at
the Club Cast3M by Gornet~\cite{Gornet:club2017}, which computes from a
meshed cross-section its area and centroid, its principal second moments, the
warping function, the torsional rigidity, the position of the torsion centre
-- located through the Maxwell--Betti reciprocal theorem -- and the shear
correction coefficients, i.e. everything a Timoshenko beam model needs. It is
validated there on triangular, rectangular, L- and C-sections, on hollow
sections with fillets, and on NACA aerofoil profiles.

That procedure is not distributed with \texttt{Cast3M} and the presentation
gives only extracts of it, so it would have to be rebuilt from the slides.
The construction above is the core of it, and rebuilding the rest is a good
project: see the exercises.

%%%
\section{Material coefficients varying with temperature}

Suppose now that a law correlates a material property $P$ with the temperature
$\theta$ -- the Young modulus $E(\theta)$, say, or the dilation coefficient
$\alpha(\theta)$. The question is how to compute the spatial distribution of
$P$ in the presence of an inhomogeneous temperature field $\theta(x)$, i.e.
\[ P(x)=P(\theta(x)) \]

\noindent The construction is the one used above for the inclusion, with the
masks replaced by an interpolation of the measured law. In practice there are
two routes: either evaluate the law directly as an algebraic expression in the
temperature field, or tabulate it as an \op{evolution} and interpolate with
\op{ipol}. The second is what you want when the law comes from measurements,
and is what the data-reading section of chapter~\ref{chap:importexport} is
for.

\framepart{Summary}

The whole chapter is one sequence of five operators, and it is the same
sequence for elasticity and for heat conduction:

\begin{center}
\begin{tabular}{@{}l@{\hspace{8mm}}l@{\hspace{8mm}}l@{}}
                     & \textbf{elasticity} & \textbf{thermal} \\[1mm]
model                & \op{mode \ldots mecanique} & \op{mode \ldots thermique} \\
material             & \op{mate 'YOUN' 'NU'}      & \op{mate 'K' 'RHO' 'C'} \\
matrix               & \op{rigi}                  & \op{cond}, \op{capa} \\
boundary conditions  & \op{bloq} + \op{depi}      & \op{bloq 'T'} + \op{depi} \\
solution             & \op{reso}                  & \op{reso} \\
\end{tabular}
\end{center}

\begin{gbox*}
\begin{itemize}
  \item \textbf{Imposed conditions add unknowns rather than removing them.}
  They are handled by Lagrange multipliers, and those multipliers are the
  reactions, recovered with \op{reac}.
  \item \textbf{Every load that is not a nodal force is built the same way.}
  Pressure, weight, thermal expansion: form a field, integrate it against
  the shape functions with the operator that belongs to it (\op{pres},
  \op{mass}, \op{thet}), and add the result to the right-hand side.
  \item \textbf{Spatially varying data is built from the coordinates.}
  Whether it is a load or a material parameter: \op{coor} and \op{masq} to
  build it, \op{chan} to move it to the Gauss points.
\end{itemize}
\end{gbox*}

\framepart{Exercises}

\exercise{Stresses and strains in a plate with a hole} Run the program
  presented in this chapter and:
  \begin{itemize}
  \item check the stress distribution for all types of applied force: body
    force, applied surface pressure, tangential traction;
  \item compute the solution for a given resultant force applied to the upper
    face of the plate, using both \op{forc[e]} and \op{pres[sion]}. Explain
    the difference between the two;
  \item compare the computed stress concentration factor with the
    theoretical value $K_t=3$, and study how it changes as the ratio
    $r/L$ of hole radius to plate width decreases.
  \end{itemize}
  See the file \texttt{traction\_trou.dgibi}.

\exercise{Mesh convergence} Repeat the computation with \op{nel} equal to
  5, 10, 20 and 40 and plot the maximum von Mises stress against the number of
  degrees of freedom. At which refinement is the result converged to within
  1\%? Compare the convergence of the displacement with that of the stress and
  explain why they differ.

\exercise{Elastic and thermal example files} Run the elasticity example
  files \texttt{elas*.dgibi} and the thermal ones \texttt{ther*.dgibi} from
  the \texttt{cast3m\_dir/dgibi} directory of your installation, and read the
  notes in \texttt{cast3m\_dir/dgibi/annotated\_testing\_files\_98.pdf} and
  on the \textbf{Exemples} page of the website.

\exercise{Torsion by the thermal analogy} Following
  section~\ref{sec:elas:torsion}, compute the warping function and the
  torsional rigidity $J$ of a \emph{circular} section, for which the exact
  answer is $J=\pi R^{4}/2$ and the warping is identically zero -- a good
  check that your boundary term is right. Then do a square section, where
  $J \simeq 0.1406\,a^{4}$ and the warping is not zero, and plot it. Finally
  a thin rectangle, and compare with $J\simeq\frac{1}{3}ab^{3}$.

  \emph{Going further:} rebuild the rest of \texttt{@Torsion} -- centroid,
  principal axes and second moments, torsion centre, shear coefficients --
  as a procedure of your own taking a meshed section and returning a table.
  It is a satisfying small project and the result is genuinely useful.

\exercise{Thermoelastic coupling} Compute the temperature field of the
  plate with the hole with the boundary conditions of
  section~\ref{sec:elas:thermal}, then use it as a loading for the elastic
  problem through a thermal dilation coefficient \texttt{'ALPH'}. Where are
  the largest stresses, and why are they not where the temperature is
  highest?
